\documentclass[12pt]{article}
\usepackage[utf8]{inputenc}
\usepackage[margin=1in]{geometry}
\usepackage{amsmath}
\usepackage{graphicx}
\usepackage{booktabs}
\usepackage{hyperref}
\usepackage[round]{natbib}

\title{Variable PC--CbN Synaptic Weights Enable Concurrent Rate and Timing Codes in Cerebellar Output}
\author{Author A$^{1}$, Author B$^{1}$, Author C$^{1,2}$ \\
\small $^{1}$Department of Neurobiology, Northwestern University, Evanston, IL \\
\small $^{2}$Northwestern University Interdepartmental Neuroscience Program}
\date{}

\begin{document}
\maketitle

\begin{abstract}
Purkinje cells (PCs) of the cerebellar cortex provide the sole output to cerebellar nuclei (CbN) neurons, yet the functional implications of the synaptic weight distribution at this critical synapse are poorly understood. Here, we report that unitary PC--CbN input sizes in juvenile mice (P23--32, $n = 83$) follow a highly skewed distribution that differs significantly from the more uniform distribution in young mice (P10--20, $n = 74$; $p < 0.0001$, Kolmogorov-Smirnov test). Using dynamic clamp with realistic PC firing statistics derived from in vivo recordings, we demonstrate that large inputs (30 nS) strongly influence both the rate and timing of CbN neuron firing, transiently eliminating output for several milliseconds after each PC spike. The PC refractory period creates a brief pre-spike disinhibition that elevates CbN firing before each inhibitory event, generating a characteristic disinhibition-then-inhibition temporal motif. Nonuniform input sizes elevate baseline CbN firing rates by increasing inhibitory conductance variability and reduce the relative influence of PC synchrony on firing rate. Nevertheless, synchronizing even two large inputs significantly increases CbN output. These findings demonstrate that the variable synaptic weight distribution at the PC--CbN synapse enables the cerebellum to support concurrent rate and timing codes, with individual large inputs exerting outsized influence on output timing.
\end{abstract}

\section{Introduction}

The cerebellum coordinates movement, balance, and motor learning through a circuit architecture in which Purkinje cells (PCs) of the cerebellar cortex provide the sole inhibitory output to glutamatergic projection neurons of the cerebellar nuclei (CbN) \citep{ito2006cerebellar}. PCs fire spontaneously at high rates (typically 30--100 Hz in vivo) and converge onto CbN neurons with estimates of approximately 30--50 PC inputs per CbN cell. A central debate in cerebellar physiology concerns whether PCs encode information primarily through firing rate modulation (rate codes) or through the synchronous timing of spikes across PC populations (timing codes), or some combination of both \citep{heck2013cerbellar, deschutter1999cerebellar}.

The rate code hypothesis posits that graded changes in PC firing rate produce graded changes in CbN neuron firing through tonic inhibition, with reduced PC activity permitting CbN rebound firing and increased PC activity suppressing CbN output \citep{gauck2000control}. The timing code hypothesis proposes that brief synchronous pauses in PC firing, even without changes in mean rate, can produce precisely timed CbN spikes via transient disinhibition \citep{person2012purkinje, bengtsson2011purkinje}. These hypotheses have been treated as largely separate mechanisms, but biological PC--CbN synapses may support both simultaneously if the synaptic weight distribution permits individual inputs to exert sufficient influence on CbN timing while populations of inputs collectively convey rate information.

Previous electrophysiological characterizations of PC--CbN inputs found them to be relatively small and uniform in size, but these studies were conducted primarily in young animals (P13--29) with limited sample sizes ($n \sim 30$) \citep{person2012purkinje, telgkamp2004maintenance}. Similarly, prior dynamic clamp studies used 40 uniform-sized PC inputs to explore CbN firing control \citep{gauck2000control}, which may not reflect the actual synaptic weight distribution in more mature animals. The functional consequences of having a heterogeneous, potentially skewed distribution of input sizes on CbN neuron firing control have not been investigated.

Here, we address this gap by (1) measuring unitary PC--CbN input sizes with substantially larger sample sizes in both young and juvenile mice, (2) using dynamic clamp with realistic PC firing statistics derived from in vivo recordings to test how inputs of different sizes influence CbN rate and timing, and (3) systematically exploring how nonuniform input sizes interact with PC synchrony to shape cerebellar output.

\section{Methods}

\subsection{Animals and Slice Preparation}

C57BL/6 wild-type mice (Charles River) of both sexes were used at two age ranges: P10--P20 (young) and P23--P32 (juvenile). Acute sagittal cerebellar slices (170 $\mu$m) were prepared in warm choline ACSF at 34$^\circ$C. All recordings were performed at 34--35$^\circ$C with a perfusion rate of 3--5 mL/min.

\subsection{Electrophysiology}

Whole-cell voltage-clamp recordings were made from large glutamatergic CbN projection neurons ($>$70 pF) in the lateral and interposed deep cerebellar nuclei. The internal solution was CsCl-based with high chloride ($E_\text{Cl} = 0$ mV), containing 110 mM CsCl, 20 mM Cs-BAPTA (to prevent plasticity), 2 mM QX-314, and 0.2 mM D600. Holding potential was $-30$ to $-40$ mV. Series resistance was compensated up to 80\%. Pharmacological blockers included NBQX (5 $\mu$M, AMPARs), (R,S)-CPP (2.5 $\mu$M, NMDARs), strychnine (1 $\mu$M, glycine receptors), and CGP 55845 (1 $\mu$M, GABA$_\text{B}$ receptors).

Unitary PC--CbN input sizes were measured using minimal stimulation of PC axons in the white matter. Stimulation intensity was adjusted to isolate individual inputs based on all-or-none IPSC responses with stable amplitude (Table~\ref{tab:inputs}).

\begin{table}[htbp]
\centering
\caption{Unitary PC--CbN input size distributions by age group.}
\label{tab:inputs}
\begin{tabular}{lccc}
\toprule
Age Group & Range & $N$ Inputs & Distribution Shape \\
\midrule
Young (P10--20) & P10--P20 & 74 & Relatively uniform \\
Juvenile (P23--32) & P23--P32 & 83 & Highly skewed \\
\multicolumn{2}{l}{K-S test: $p < 0.0001$} & & \\
\bottomrule
\end{tabular}
\end{table}

\subsection{Dynamic Clamp}

Dynamic clamp experiments (P26--32) used a configuration designed to approximate the juvenile input distribution: 16 small inputs (3 nS each), 10 medium inputs (10 nS each), and 2 large inputs (30 nS each), for a total of 28 inputs \citep{sharp2011dynamic}. Firing statistics for each input were derived from 10 PCs recorded in vivo, with interspike interval distributions fit by lognormal functions.

Rate coding was tested by varying the firing rate of one input size class (small, medium, or large) while holding others at 80 spikes/s. Synchrony was manipulated by introducing correlated pauses across inputs at varying fractions (0--100\%). Cross-correlograms between individual PC input spike trains and CbN neuron output were computed to assess timing influence.

\subsection{Computational Modeling}

A simple integrate-and-fire CbN neuron model with conductance-based inhibitory inputs was used to complement the dynamic clamp experiments. The model faithfully reproduced experimental findings and allowed systematic exploration of input number, size distribution, synchrony level, and firing rate.

\section{Results}

\subsection{Juvenile PC--CbN Inputs Follow a Highly Skewed Distribution}

Minimal stimulation recordings revealed that unitary PC--CbN input sizes in juvenile mice (P23--32) follow a markedly different distribution from those in young mice (P10--20). The young distribution was relatively narrow with modest variability, consistent with prior reports \citep{person2012purkinje}. In contrast, the juvenile distribution was highly skewed, with many medium inputs and a long tail toward large inputs (Kolmogorov-Smirnov test, $p < 0.0001$). In a single P27 cell, inputs ranged from 720 pA (small) through 1690 pA (medium) to 6280 pA (large), illustrating the range within a single CbN neuron.

\subsection{PC Refractory Period Creates a Disinhibition-Inhibition Motif}

Analysis of in vivo PC firing statistics revealed that the refractory period following each PC spike creates a temporal window of reduced spike probability. This refractory period causes a brief dip in the total inhibitory conductance impinging on a CbN neuron before the next PC spike arrives. In dynamic clamp experiments, this translated to a brief elevation of CbN firing probability (disinhibition) preceding the inhibitory suppression caused by the next PC spike, generating a characteristic disinhibition-then-inhibition temporal motif. This effect was absent when artificial Poisson inputs (lacking a refractory period) were used, confirming that it arises specifically from the temporal statistics of real PC firing (Table~\ref{tab:refractory}).

\begin{table}[htbp]
\centering
\caption{Effect of PC refractory period on CbN neuron response.}
\label{tab:refractory}
\begin{tabular}{lcccc}
\toprule
Input Type & Pre-spike Conductance & Post-spike Effect & CbN Motif \\
\midrule
Real PC (with refractory) & Dip (disinhibition) & Inhibition & Elevation then suppression \\
Poisson (no refractory) & Flat & Inhibition & Suppression only \\
\bottomrule
\end{tabular}
\end{table}

\subsection{Conductance Variability Elevates Baseline CbN Firing}

A key finding was that nonuniform input sizes elevate baseline CbN firing rates even when the mean inhibitory conductance is held constant. When different input configurations all producing a mean conductance of 56 nS were compared, CbN neurons fired at the highest rates with fewer, larger inputs (high variability) and at the lowest rates with constant conductance (zero variability). This occurs because conductance fluctuations create brief windows of reduced inhibition during which CbN neurons can fire, an effect that increases with the amplitude of fluctuations ($n = 7$ cells).

\subsection{Rate Code Transmission Scales with Input Size}

Varying the firing rate of one input size class while holding others at 80 spikes/s revealed that all input sizes can transmit rate codes---changes in PC firing rate inversely modulated CbN output. However, the sensitivity scaled dramatically with input size (Table~\ref{tab:rate}). Large inputs (2 $\times$ 30 nS) produced the strongest modulation of both mean inhibitory conductance and CbN firing rate, despite representing only 2 of 28 total inputs. Small inputs (16 $\times$ 3 nS) produced modest modulation, and medium inputs (10 $\times$ 10 nS) produced intermediate effects.

\begin{table}[htbp]
\centering
\caption{Rate code transmission by input size class.}
\label{tab:rate}
\begin{tabular}{lccc}
\toprule
Input Class & $N \times$ Conductance & Conductance Modulation & CbN Rate Sensitivity \\
\midrule
Small & $16 \times 3$ nS & Small & Low \\
Medium & $10 \times 10$ nS & Moderate & Moderate \\
Large & $2 \times 30$ nS & Large & High \\
\bottomrule
\end{tabular}
\end{table}

\subsection{Cross-Correlograms Reveal Size-Dependent Timing Influence}

Cross-correlograms between individual PC input spike trains and CbN output spikes revealed that the timing influence of a single input scales with its size. Large inputs (30 nS) produced strong modulation with a prominent disinhibition peak followed by deep suppression lasting several milliseconds, effectively eliminating CbN firing for that period. Medium inputs (10 nS) showed moderate modulation, and small inputs (3 nS) produced weak, brief effects. The refractory-period-driven disinhibition peak was visible for all sizes but most prominent for large inputs.

\subsection{Nonuniform Inputs Reduce Relative Synchrony Influence}

When all 40 inputs were uniform (5 nS each), increasing synchrony from 0\% to 100\% produced a monotonic increase in both the coefficient of variation (CV) of inhibitory conductance and CbN firing rate. With nonuniform inputs (16 small + 10 medium + 2 large), the baseline CV was already elevated due to the inherent variability of the input sizes, and the relative effect of synchrony on firing rate was smaller. However, this did not mean that synchrony was ineffective: synchronizing even two large inputs produced a measurable increase in CbN firing, demonstrating that timing-based coding does not require widespread PC synchrony when large synaptic inputs are present \citep{shin2007dynamic}.

\section{Discussion}

This study reveals that the PC--CbN synapse undergoes a developmental transformation from a relatively uniform to a highly skewed weight distribution, and that this heterogeneity has profound consequences for how the cerebellum encodes information through both rate and timing codes.

\subsection{Developmental Emergence of Synaptic Weight Heterogeneity}

The discovery that juvenile PC--CbN inputs follow a highly skewed distribution, significantly different from the more uniform distribution in young animals, has important implications for interpreting prior studies. The earlier characterizations using young animals and smaller sample sizes naturally underestimated the variability present in the mature circuit \citep{person2012purkinje, telgkamp2004maintenance}. The skewed distribution, with a long tail toward large inputs, means that a small number of PC synapses can exert disproportionate control over CbN neuron firing---a feature that is functionally distinct from the uniform-input assumption used in previous modeling and dynamic clamp studies \citep{gauck2000control, abbasi2017robust}.

\subsection{Concurrent Rate and Timing Codes}

Our results demonstrate that the variable weight distribution enables simultaneous support of rate and timing codes at the PC--CbN synapse \citep{heck2013cerbellar, deschutter1999cerebellar}. Rate codes are supported because changes in firing rate of any input size class modulate CbN output, with large inputs having outsized influence. Timing codes are supported because individual large inputs can precisely control CbN spike timing through the disinhibition-inhibition motif, and synchronizing even a small number of large inputs produces significant effects on CbN firing \citep{person2012purkinje, bengtsson2011purkinje}.

The disinhibition-inhibition motif arising from the PC refractory period adds a previously unrecognized dimension to timing code mechanisms. Rather than simply suppressing CbN firing, each PC spike first creates a brief window of elevated firing probability before suppression, producing precisely timed CbN output locked to the PC firing pattern. This mechanism operates independently of synchrony across multiple PCs and is most pronounced for large inputs, providing a channel for single-PC timing information to be transmitted to downstream targets \citep{raman2000resurgent}.

\subsection{Implications for Cerebellar Computation}

The elevation of baseline CbN firing by conductance variability, even at matched mean inhibitory conductance, suggests that the synaptic weight distribution itself is a computational parameter of the cerebellar circuit. A more heterogeneous distribution produces higher baseline firing rates, effectively increasing the dynamic range within which rate codes can operate. Simultaneously, the large inputs provide temporal precision for timing-based computations. This dual functionality may be particularly important for motor coordination tasks that require both graded force control (rate coding) and precise temporal sequencing (timing coding) \citep{ito2006cerebellar, wu2019bhatt}.

\subsection{Limitations}

Our dynamic clamp approximation used 28 inputs (16 small, 10 medium, 2 large), which is fewer than the estimated 30--50 PCs converging onto a single CbN neuron in vivo. The in vivo PC statistics were derived from 10 cells, which may not capture the full range of PC firing behaviors. The slice preparation removes ongoing excitatory and neuromodulatory inputs that would influence CbN firing in intact circuits.

\section{Conclusion}

The highly skewed distribution of PC--CbN synaptic weights in juvenile mice enables the cerebellar output to simultaneously encode rate and timing information. Large inputs dominate timing control through a disinhibition-inhibition motif driven by the PC refractory period, while the collective activity of inputs of all sizes transmits rate codes. Conductance variability from nonuniform inputs elevates baseline CbN firing, and synchronizing even a few large inputs significantly modulates output. These findings unify the rate and timing code hypotheses of cerebellar function and identify the synaptic weight distribution as a key computational feature of the cerebellar circuit.

\bibliographystyle{plainnat}
\bibliography{references}

\end{document}
